% status: FULL
% Authoritative LaTeX; edit this file, not the historical Markdown.
\chapter{Embeddings and Normalization}\label{chapter-7-embeddings-and-normalization}\label{app:embeddings-norm}
\Appref{tokens} ended with integers. A tokenizer transformed UTF-8 bytes into
token identities under a model-specific vocabulary contract. \Appref{smallest-model}
used one of those identities to select a hidden vector, but its purpose
was to make an entire language-model forward pass small enough to
calculate by hand. Now we can inspect that first numerical step as an
engine operation in its own right.

A token ID is not a small activation. Multiplying the integer
\texttt{128} by a model weight does not recover what token 128
``means.'' The integer is an addressable symbol. The model's
\textbf{embedding table} stores one learned floating-point row for each
symbol. A checked lookup crosses from discrete token space into the
continuous model space in which the network computes.

Once there, scale matters. A vector and a hundred-times-larger copy
point in the same direction, but their components drive later arithmetic
at very different magnitudes. Transformer architectures commonly place
normalization around learned transformations to control that scale. This
chapter builds \textbf{RMSNorm} from squares, a mean, a square root, a
reciprocal, and a learned element-wise scale. There is no hidden
framework call.

These operators are short. Their contracts are not. We need to know
which axes mean vocabulary and model width, which bytes a row uses,
whether the result aliases weights, what happens to an empty dimension,
where epsilon appears, which precision performs the reduction, and what
finite values can still overflow. Those decisions become assumptions of
every later Transformer chapter.

\begin{quote}
\textbf{FIRST PRINCIPLE} Tokenization chooses a discrete model symbol.
Embedding lookup materializes that symbol's learned model-space vector.
Normalization then changes the vector's scale under an explicit
numerical contract; it does not invent its meaning or clip its
coordinates.
\end{quote}

\section{The first learned vector}\label{the-first-learned-vector}

\EngineFigure{ch07-journey}{Token identity, copied embedding and normalized activation occupy different boundaries.}{fig:ch07-journey}

The visual edition follows one deliberately small table through the
whole chapter. It has three vocabulary rows and four coordinates per
row; token 1 selects the existing RMSNorm hand vector
\texttt{{[}1,-2,3,-4{]}}. This is a new connecting fixture, not a
replacement for \Appref{smallest-model}'s three-dimensional tiny model or the original
\Appref{embeddings-norm} table tests. Keeping these identities explicit prevents a
polished diagram from silently changing the engine being explained.

Read the plate as three contracts. An integer identifies a row. Lookup
copies that row into activation storage. RMSNorm creates another
activation under a specified equation. The dashed next-step arrow is a
curriculum boundary, not an implemented Q/K/V call. This is the
model-semantic view; the ownership and physical-address views below
explain what its short arrows cost.

The complete boundary is visible in the
\href{https://github.com/hermonai/inside-the-llm-engine/blob/astra-visual-rewrite/diagrams/transformer/token-to-model-space.txt}{token-to-model-space
diagram}:

The token identity selects a parameter row. The copied row becomes an activation. Normalization transforms that activation while leaving the embedding table unchanged.

The tokenizer owns text segmentation and the mapping from pieces to
integer IDs. The numerical model owns the embedding table. An ID has
meaning only relative to that pair of artifacts: the same integer can
select an unrelated row in a different model.

Let \(V_{\mathrm{vocab}}\) be the vocabulary size and \(D\) be the model
or embedding dimension. The table is a learned parameter matrix

\begin{equation}
\mathbf{E}\in\mathbb{R}^{V_{\mathrm{vocab}}\times D}.
\end{equation}

For a valid token ID \(t\), lookup produces

\begin{equation}
\mathbf{x}=\mathbf{E}_{t,:}\in\mathbb{R}^{D},
\qquad 0\le t<V_{\mathrm{vocab}}.
\end{equation}

The colon means every column in row \(t\). Written element by element,

\begin{equation}
x_j=E_{t,j},
\qquad 0\le j<D.
\end{equation}

This equation specifies selection, not arithmetic over every table row.
A one-hot vector times \(\mathbf{E}\) is a mathematically equivalent
description, but constructing that mostly-zero vector and performing a
dense product would be a perverse implementation. The engine already has
the row index.

The output is the first \textbf{residual activation}. In a decoder, the
\textbf{residual stream} is the persistent model-width state carried
from one Transformer sublayer to the next. \Appref{embeddings-norm} establishes its
width and normalization; it does not yet build any of the
transformations that consume it.

\section{What is physically stored?}\label{what-is-physically-stored}

The table contains floating-point parameters learned during training. It
does not contain source strings, dictionary definitions, or
human-readable facts. Rows can acquire useful geometric relationships,
but an inference engine needs no story about their semantics to execute
lookup. It needs a tensor with valid metadata and a token in range.

The
\href{https://github.com/hermonai/inside-the-llm-engine/blob/astra-visual-rewrite/diagrams/transformer/embedding-logical-layout.txt}{logical
layout} separates the two axes:

\begin{itemize}
\tightlist
\item
  axis 0 has length \(V_{\mathrm{vocab}}\) and is selected by token
  identity;
\item
  axis 1 has length \(D\) and becomes the model-space vector.
\end{itemize}

That is why the semantic shape is \texttt{{[}V,D{]}}, rather than
\texttt{{[}D,V{]}}. The convention matches the operator: selecting
\texttt{table{[}t,:{]}} returns \texttt{D} values. Another system can
store axes differently, but then its metadata and kernel must make that
orientation explicit.

In Tensor Substrate v1, an \texttt{OwnedTensor} uses canonical row-major
storage. An embedding owner therefore reports

\begin{verbatim}
shape   = [V,D]
strides = [D,1]        // measured in elements
dtype   = f32
\end{verbatim}

For base offset zero, the canonical physical element offset is

\begin{equation}
o(t,j)=tD+j.
\end{equation}

Because one \texttt{f32} occupies four bytes, the corresponding byte
displacement from the allocation's first element is

\begin{equation}
b(t,j)=4(tD+j)\quad[\mathrm{bytes}].
\end{equation}

Those equations describe the canonical case. The reference operator
accepts a general validated \texttt{TensorView}, so its actual logical
mapping remains \Appref{tensors}'s stride equation:

\begin{equation}
o(t,j)=b_0+t s_0+j s_1,
\end{equation}

where \(b_0\) is the view's base element offset and \(s_0,s_1\) are
nonnegative element strides. The
\href{https://github.com/hermonai/inside-the-llm-engine/blob/astra-visual-rewrite/diagrams/transformer/embedding-physical-layout.txt}{physical-layout
diagram} puts the canonical shortcut beside the general rule. Code calls
checked \texttt{get2(t,j)} instead of reconstructing either formula in
the operator. That keeps extent validation and logical addressing in the
tensor substrate.

For a canonical table, one selected row is physically contiguous:
reading columns \texttt{0..D} walks adjacent \texttt{f32} values. A
strided table can place padding or another logical arrangement between
columns. A zero-stride view can repeat one storage value across an axis.
Tensor Substrate v1 admits those immutable views, so the broad reference
operator defines their behavior instead of quietly assuming stride one.

\section{The checked lookup
contract}\label{the-checked-lookup-contract}

\EngineFigure{ch07-lookup}{Token 1 selects four F32 values at element offsets 4 through 7 and byte displacements 16 through 28.}{fig:ch07-lookup}

The table and activation must not share a label merely because both are
matrices of numbers. The learned table is \texttt{{[}V,D{]}}; a sequence
of selected activations is \texttt{{[}T,D{]}}. Here one token selects
\texttt{{[}D{]}}. In the plate, the ochre cells remain parameters and
the green cells are a new owner. The address labels are displacements
within the canonical table, not allocator addresses and not offsets
within the returned vector, whose own first element is zero.

\texttt{\seqsplit{embedding\_lookup\_reference}} accepts an immutable
table view and a \texttt{TokenId}.

\Needspace{12\baselineskip}

It requires:

\begin{enumerate}
\def\labelenumi{\arabic{enumi}.}
\tightlist
\item
  table rank exactly two;
\item
  vocabulary size \(V_{\mathrm{vocab}}>0\);
\item
  model dimension \(D>0\);
\item
  token ID in the half-open interval \([0,V_{\mathrm{vocab}})\);
\item
  a view whose reachable extent was already validated by Tensor
  Substrate v1.
\end{enumerate}

It returns a canonical \texttt{OwnedTensor} of shape \texttt{{[}D{]}}.
Wrong rank, empty axes, and an invalid ID produce structured
\texttt{EmbeddingError} variants. A token equal to
\(V_{\mathrm{vocab}}\) is already out of range; the last valid ID is
\(V_{\mathrm{vocab}}-1\). Validation precedes element reads, so ordinary
bad input cannot become a panic or a speculative out-of-bounds access.

The core loop is deliberately unsurprising:

\begin{verbatim}
let mut output = Vec::with_capacity(model_dimension);
for column in 0..model_dimension {
    output.push(*table.get2(row, column)?);
}
OwnedTensor::from_vec(vec![model_dimension], output)
\end{verbatim}

The important code is also around the loop: rank and dimension checks,
checked conversion from \texttt{TokenId}, output shape construction, and
the owned return type. A five-line numerical body does not imply a
five-line operator contract.

\begin{quote}
\textbf{BUILD IT} Use
\href{https://github.com/hermonai/inside-the-llm-engine/blob/astra-visual-rewrite/labs/lab-30-inspect-embedding-row.md}{Lab
30} to calculate canonical and strided offsets, then
\href{https://github.com/hermonai/inside-the-llm-engine/blob/astra-visual-rewrite/labs/lab-31-checked-embedding-lookup.md}{Lab
31} to audit every validation boundary. An off-by-one comparison must
fail at token \texttt{V}, not during a later storage access.
\end{quote}

\section{View or copy?}\label{view-or-copy}

The table already owns the selected values. Why allocate another vector?

One possible interface returns a
\texttt{\seqsplit{TensorView<'weights>}} of the row. That avoids copying
\(D\) elements. It also aliases model parameters and carries a lifetime
tied to their owner. An immutable consumer could use it safely, but a
residual stream is not just an immutable observation: later sublayers
produce additions and other request-local state. If code expects to
mutate the row as an activation, aliasing model weights would be
catastrophic.

The selected \Appref{embeddings-norm} policy is therefore explicit materialization:

For each input position $t$, gather $X_{t,:}=E_{x_t,:}$. Repeated IDs give equal row values but distinct output positions.

The fresh owner costs one read and one write per element. In the
simplest cold-payload accounting, a canonical \texttt{f32} lookup moves
approximately

\begin{equation}
Q_{\mathrm{lookup}}\approx 8D\quad[\mathrm{bytes}],
\end{equation}

counting \(4D\) input bytes and \(4D\) output bytes. This is not
measured memory traffic: cache lines, write allocation, allocator
metadata, and warm data can change physical transfers. It is a useful
statement of the policy's payload cost.

In exchange, downstream mutation cannot alter \(\mathbf{E}\), the
activation can outlive the lookup borrow, and request cleanup has one
obvious owner. The
\href{https://github.com/hermonai/inside-the-llm-engine/blob/astra-visual-rewrite/diagrams/transformer/embedding-view-vs-copy.txt}{view-versus-copy
diagram} records both options rather than presenting the copy as free.

A database-table analogy can help for one moment: the embedding table is
long-lived storage, and an ID chooses a row. The analogy stops there.
These are learned numerical weights, not relational records; there is no
SQL predicate, transaction, schema evolution, or human field meaning.
The actual mechanism remains tensor indexing and copying.

\begin{quote}
\textbf{PROVE IT}
\href{https://github.com/hermonai/inside-the-llm-engine/blob/astra-visual-rewrite/labs/lab-32-embedding-view-vs-copy.md}{Lab
32} mutates the returned \texttt{OwnedTensor} and verifies that the
original table is unchanged. Ownership is tested behavior, not a
comment.
\end{quote}

\section{From one token to a
sequence}\label{from-one-token-to-a-sequence}

\EngineFigure{ch07-sequence}{Repeated token identities produce equal rows at distinct positions in one activation allocation.}{fig:ch07-sequence}

Positions and identities are separate axes of reasoning. Positions 0 and
2 can contain the same token ID without sharing mutable output cells. In
this fixture, changing \texttt{X{[}0,0{]}} cannot change
\texttt{X{[}2,0{]}}, even though lookup originally wrote the same value
into both. An optimization that deduplicated parameter reads would still
owe the same logical output and ownership contract.

Inference begins with more than one prompt token even though later
generation produces one new token at a time. The single-row rule
generalizes without introducing attention.

For the token sequence

\begin{equation}
\mathbf{t}=[t_0,t_1,\ldots,t_{T-1}],
\end{equation}

sequence embedding produces

\begin{equation}
\mathbf{X}\in\mathbb{R}^{T\times D},
\qquad X_{ij}=E_{t_i,j},
\end{equation}

for \(0\le i<T\) and \(0\le j<D\). The shapes are:

\begin{verbatim}
token identities   [T]
embedding table    [V,D]
result             [T,D]
\end{verbatim}

\texttt{\seqsplit{embedding\_sequence\_reference}} validates the same
positive table dimensions, checks prospective output element count
\(TD\), then checks each token before copying its logical row. It
preserves order and repetition. Tokens \texttt{{[}3,0,3{]}} produce rows
\texttt{{[}E3,E0,E3{]}}; it does not deduplicate them or sort for
locality.

An empty token sequence has \(T=0\) and returns an owned tensor with
shape \texttt{{[}0,D{]}}. That result has a well-defined shape and no
values. This does not make an empty model vocabulary valid: \(V=0\)
leaves no token identity selectable, and \(D=0\) would not create a
model-space vector. The axes have independent contracts.

The implementation is intentionally only a gather into
\texttt{{[}T,D{]}}. It does not create batches, positions, masks, Q/K/V,
or a Transformer block. Those concepts need separate semantic and
ownership work.

\section{The residual-stream
invariant}\label{the-residual-stream-invariant}

\EngineFigure{ch07-owners}{Model parameters, immutable views and request-owned activation tensors have distinct lifetimes.}{fig:ch07-owners}

This diagram uses the real storage types, \texttt{OwnedTensor} and
\texttt{TensorView}, not an imagined Transformer class hierarchy. The
dashed edge points from a view to the owner it borrows. Solid edges
represent a copy or numerical transformation. Neither means transferring
ownership of model weights into a request. Model storage can serve many
requests; request cleanup must release only its own activation owners.

Once lookup produces \(\mathbf{x}_0\in\mathbb{R}^{D}\) for a token,
\(D\) becomes one of the decoder's central dimensional invariants.
Transformer sublayers can create internal tensors with other shapes, but
a value added back into the residual stream must return to model width
\(D\).

The
\href{https://github.com/hermonai/inside-the-llm-engine/blob/astra-visual-rewrite/diagrams/transformer/residual-stream-width.txt}{residual-width
diagram} states that bounded claim. ``Persistent'' means the state is
carried through the layer sequence for this request, not that it has
model lifetime or durable storage. It is activation data. The embedding
table and normalization weights are parameters; the residual vectors are
request-owned values.

This distinction is summarized in the
\href{https://github.com/hermonai/inside-the-llm-engine/blob/astra-visual-rewrite/diagrams/transformer/parameters-vs-activations.txt}{parameter/activation
diagram}:

Read-only parameter storage outlives individual calls; each request creates and releases its own activation owners. Figure~\ref{fig:parameter-owners} gives the corresponding ownership relation.

Inference does not update the learned parameters. Multiple requests may
read them concurrently. Each request owns the activations it creates and
eventually releases them. Later chapters will give that distinction
consequences for loading, placement, caches, and scheduling; here it
already prevents parameter aliasing.

\section{Why scale needs a
contract}\label{why-scale-needs-a-contract}

Consider two vectors:

\begin{equation}
\mathbf{x}_1=[1,2,-1,2],
\qquad
\mathbf{x}_2=[10,20,-10,20].
\end{equation}

The second is \(10\mathbf{x}_1\). Their coordinate pattern and direction
agree, but every component of the second is ten times larger. Repeated
learned transformations and residual additions can produce activation
scales that vary across layers and inputs. Later matrix products
multiply those values by many weights, so scale affects representable
range and numerical behavior.

Normalization provides a defined rescaling based on a vector statistic.
It is not clipping. A clipping operator might independently replace
every value above a threshold with that threshold, changing coordinates
selectively. RMSNorm computes one scalar from the entire vector, scales
every coordinate by that shared value, and then applies learned
per-coordinate weights. Large components remain large relative to small
components unless learned scale changes that relation.

The original RMSNorm paper by Biao Zhang and Rico Sennrich proposed
removing LayerNorm's recentering step while retaining RMS-based
rescaling. That history is useful, but the inference engine still needs
the exact forward equation.

\section{Root mean square from first
principles}\label{root-mean-square-from-first-principles}
\index{root mean square}

Let

\begin{equation}
\mathbf{x}=[x_0,x_1,\ldots,x_{D-1}]\in\mathbb{R}^{D},
\qquad D>0.
\end{equation}

First square each component. Squaring makes every contribution
nonnegative, so opposite signs cannot cancel. Next take the arithmetic
mean:

\begin{equation}
m_2(\mathbf{x})=\frac{1}{D}\sum_{i=0}^{D-1}x_i^2.
\end{equation}

The quantity \(m_2\) has squared units relative to \(\mathbf{x}\).
Taking the square root restores the original units:

\begin{equation}
\operatorname{RMS}(\mathbf{x})
=\sqrt{\frac{1}{D}\sum_{i=0}^{D-1}x_i^2}.
\end{equation}

That sequence---square, mean, root---is the name \textbf{root mean
square}. The
\href{https://github.com/hermonai/inside-the-llm-engine/blob/astra-visual-rewrite/diagrams/transformer/rms-calculation-pipeline.txt}{RMS
pipeline} makes each lowering step visible.

Why take a mean rather than only a sum? If a vector repeated the same
magnitude at every coordinate, the sum of squares would grow with \(D\).
Dividing by \(D\) makes its RMS equal that coordinate magnitude,
independent of vector length. For \texttt{{[}4,4,4,4{]}}, RMS is 4, not
8.

RMS alone has a zero-denominator problem. For the zero vector, every
square and their mean are zero. The reciprocal required for
normalization would be undefined. RMSNorm therefore defines an
epsilon-stabilized denominator. In this book and the checked operator,
epsilon is \emph{inside} the square root:

\begin{equation}
r=\frac{1}{\sqrt{m_2(\mathbf{x})+\epsilon}},
\qquad \epsilon>0.
\end{equation}

The expression

\begin{equation}
\frac{1}{\sqrt{m_2(\mathbf{x})}+\epsilon}
\end{equation}

is a different function. It cannot be substituted just because both
contain a square root and epsilon. PyTorch's official RMSNorm semantics
and the inspected Hermon/llama.cpp paths use the first convention.

\section{The RMSNorm operator}\label{the-rmsnorm-operator}
\index{RMSNorm}

RMS rescaling alone would apply scalar \(r\) to every coordinate.
RMSNorm also has a learned scale vector

\begin{equation}
\mathbf{w}\in\mathbb{R}^{D}.
\end{equation}

The complete operator is

\begin{equation}
\boxed{
y_i=x_i
\left(
\frac{1}{\sqrt{\frac{1}{D}\sum_{j=0}^{D-1}x_j^2+\epsilon}}
\right)
w_i
},
\qquad 0\le i<D.
\end{equation}

Equivalently, if \(r\) names the scalar reciprocal RMS,

\begin{equation}
\mathbf{y}=r(\mathbf{x}\odot\mathbf{w}),
\qquad
\mathbf{x},\mathbf{w},\mathbf{y}\in\mathbb{R}^{D}.
\end{equation}

The operator has two different kinds of multiplication. The reciprocal
RMS \(r\) is one scalar derived from all of \(\mathbf{x}\). The learned
scale \(\mathbf{w}\) is element-wise: coordinate \(i\) uses \(w_i\). The
symbol \(\odot\) denotes that element-wise product, not a dot product.

Learned scale lets the trained model adjust normalized coordinates
instead of forcing every feature to use identical output scale. During
inference, \(\mathbf{w}\) is immutable parameter data. No learning
occurs inside this function.

\Needspace{28\baselineskip}

\subsection{A complete calculation by
hand}\label{a-complete-calculation-by-hand}

Use

\begin{equation}
\mathbf{x}=[1,-2,3,-4],\qquad
\mathbf{w}=[1,0.5,2,-1],\qquad
\epsilon=10^{-5}.
\end{equation}

The squares sum to

\begin{equation}
1^2+(-2)^2+3^2+(-4)^2=1+4+9+16=30.
\end{equation}

With \(D=4\),

\begin{equation}
m_2=\frac{30}{4}=7.5,
\end{equation}

and

\begin{equation}
r=\frac{1}{\sqrt{7.5+10^{-5}}}
\approx0.36514813.
\end{equation}

Applying both scales gives

\begin{equation}
\begin{aligned}
y_0&= 1(r)(1)       &&\approx 0.36514813,\\
y_1&=-2(r)(0.5)     &&\approx-0.36514813,\\
y_2&= 3(r)(2)       &&\approx 2.1908888,\\
y_3&=-4(r)(-1)      &&\approx 1.4605925.
\end{aligned}
\end{equation}

The independent Python oracle and a focused Rust test calculate the
entire vector. The printed decimals are evidence from those programs,
not unchecked mental arithmetic.

\begin{quote}
\textbf{BUILD IT} Complete
\href{https://github.com/hermonai/inside-the-llm-engine/blob/astra-visual-rewrite/labs/lab-33-compute-rms-by-hand.md}{Lab
33}. Keep exact symbolic steps separate from approximate decimal output,
then use
\href{https://github.com/hermonai/inside-the-llm-engine/blob/astra-visual-rewrite/labs/lab-34-implement-rmsnorm.md}{Lab
34} to trace those same stages through the reference operator.
\end{quote}

\section{Equation to two loops}\label{equation-to-two-loops}

\EngineFigure{ch07-passes}{Four squares reduce to a shared reciprocal RMS before the second pass writes output coordinates.}{fig:ch07-passes}

The
\href{https://github.com/hermonai/inside-the-llm-engine/blob/astra-visual-rewrite/figures/generated/ch07-passes.html}{four-step
sequence} highlights the same stages, with keyboard controls and no
autoplay. Its static plate is complete: every intermediate needed for
the hand calculation remains visible in print. The highlighted stage
changes only the explanation, not the reduction order. Partial sums
\texttt{{[}1,5,14,30{]}} correspond to increasing logical index, exactly
as the reference loop below.

The direct lowering naturally has two logical passes:

\begin{verbatim}
sum_squares = 0
for i in 0..D:
    sum_squares += x[i] * x[i]

mean_square = sum_squares / D
inverse_rms = 1 / sqrt(mean_square + epsilon)

for i in 0..D:
    y[i] = x[i] * inverse_rms * weight[i]
\end{verbatim}

Pass one reduces \(D\) values to the scalar \(r\). Pass two broadcasts
that scalar while reading the input again and reading learned weights.
The output for coordinate zero cannot be finalized before pass one has
seen the last input, because the last square contributes to its
denominator.

The
\href{https://github.com/hermonai/inside-the-llm-engine/blob/astra-visual-rewrite/diagrams/transformer/rmsnorm-two-pass.txt}{two-pass
diagram} follows the data, and the
\href{https://github.com/hermonai/inside-the-llm-engine/blob/astra-visual-rewrite/diagrams/transformer/equation-to-loop.txt}{equation-to-loop
diagram} connects symbols, tensor shapes, checked logical indices, and
execution. The reference implementation follows it closely:

\Needspace{20\baselineskip}

\begin{verbatim}
let mut sum_squares = 0.0_f32;
for index in 0..dimension {
    let value = *input.get(&[index])?;
    let square = value * value;
    sum_squares += square;
}

let mean_square = sum_squares / dimension as f32;
let inverse_rms = 1.0_f32 / (mean_square + epsilon).sqrt();

for index in 0..dimension {
    output.push(*input.get(&[index])? * inverse_rms * *weight.get(&[index])?);
}
\end{verbatim}

The source contains finite-value checks omitted from this condensed
excerpt. Those checks make the simple algorithm fail explicitly when its
arithmetic leaves the declared domain.

\section{Shape, layout, and ownership
contracts}\label{shape-layout-and-ownership-contracts}

\texttt{\seqsplit{rms\_norm\_reference}} accepts two immutable
\texttt{TensorView} values and a scalar epsilon. Its complete
preconditions are:

\begin{itemize}
\tightlist
\item
  input rank is one;
\item
  learned-scale rank is one;
\item
  their lengths agree;
\item
  \(D>0\);
\item
  epsilon is finite and greater than zero;
\item
  every logical input and weight value is finite.
\end{itemize}

The function accepts any valid nonnegative-stride rank-one view. A
stride-two input reads every other physical element. A zero-stride input
repeats one physical value \(D\) times; a zero-stride weight repeats one
learned scale. That behavior is safe because views are immutable and the
tensor constructor has proved their reachable extent. The operator does
not call \texttt{to\_contiguous} and does not disguise a layout repair
as numerical work.

The output is a new canonical \texttt{{[}D{]}} \texttt{OwnedTensor}. It
aliases neither input nor weight. That policy matches embedding output:
parameter owners are read-only, and new activation data has
request-local ownership.

Wrong rank produces \texttt{RankMismatch} naming the operand. Unequal
lengths produce \texttt{LengthMismatch} with both sizes. Empty dimension
and invalid epsilon have dedicated variants. \texttt{TensorError}
remains wrapped for substrate failures. There are also structured
numerical errors for non-finite input, square, accumulated reduction,
inverse RMS, and output.

Ordinary invalid operator inputs do not panic. Resource exhaustion
remains a different boundary: after checked element count, a real
allocator failure is not converted into a tensor-shape error by this
small teaching API.

\section{Epsilon is semantic
metadata}\label{epsilon-is-semantic-metadata}

\EngineFigure{ch07-epsilon}{Positive epsilon defines zero-input behavior but breaks exact scale invariance for tiny signals.}{fig:ch07-epsilon}

For \(\mathbf{x}=\mathbf{0}\) and positive finite epsilon,

\begin{equation}
r=\frac{1}{\sqrt{\epsilon}}
\end{equation}

is finite. Each output still equals zero because \(x_i=0\). The
\href{https://github.com/hermonai/inside-the-llm-engine/blob/astra-visual-rewrite/diagrams/transformer/epsilon-zero-vector.txt}{epsilon
diagram} shows the valid and rejected paths.

The \Appref{embeddings-norm} API rejects:

\begin{itemize}
\tightlist
\item
  \texttt{epsilon\ ==\ 0}, because the zero vector would have an
  undefined reciprocal;
\item
  \texttt{epsilon\ \textless{}\ 0}, because the square-root input can
  become negative;
\item
  NaN, because it poisons comparisons and arithmetic;
\item
  positive or negative infinity, because the operator requires finite
  configuration.
\end{itemize}

Why not accept zero for nonzero vectors? An operator contract must cover
every valid input, not only the vector a caller happens to provide
today. Positive epsilon makes zero-vector behavior defined. Requiring
positivity at entry is more predictable than conditionally accepting a
configuration based on input values.

The number itself is model configuration, not a universal constant. The
educational examples use \texttt{1e-5}; current model families can
specify other values. PyTorch even defines a dtype-based default when
its \texttt{eps} argument is omitted. An inference runtime loading
trained weights must use the model's declared convention and value.
Hermon's paged preview reads
\texttt{\seqsplit{llama.attention.layer\_norm\_rms\_epsilon}} metadata,
defaults only under its current code path, and rejects nonpositive or
non-finite results at model construction. That is validation at a
different lifecycle stage, not a license to ignore epsilon.

\begin{quote}
\textbf{ENGINEERING FAILURE} Treating epsilon as an arbitrary ``small
number'' can create a numerically different model. Moving it outside the
square root, changing it during a port, or accepting NaN metadata
violates the operator contract even if many ordinary vectors produce
superficially similar output.
\end{quote}

Use
\href{https://github.com/hermonai/inside-the-llm-engine/blob/astra-visual-rewrite/labs/lab-35-break-epsilon.md}{Lab
35} to make every invalid case and the zero-vector result executable.

\section{RMSNorm is not LayerNorm}\label{rmsnorm-is-not-layernorm}

\EngineFigure{ch07-contrast}{A uniform vector is preserved near unit scale by RMSNorm but becomes zero after LayerNorm centering with zero bias.}{fig:ch07-contrast}

The uniform-vector counterexample is more useful than a vague family
resemblance. With unit gain and zero bias, LayerNorm removes all four
equal coordinates by subtracting their mean. RMSNorm does not. After RMS
rescaling, learned gain introduces another distinction: our mixed-sign
fixture's output has RMS about \texttt{1.342}. Coordinate 3 changes sign
because its gain is negative. There is no final operation that forces
the gained vector back to unit RMS.

Readers will encounter both names. They are related normalization
operators, not spelling variants.

LayerNorm computes a mean

\begin{equation}
\mu=\frac{1}{D}\sum_{i=0}^{D-1}x_i,
\end{equation}

centers the vector with \(x_i-\mu\), and scales by a variance-like
statistic before learned affine parameters. RMSNorm omits that
recentering step. It uses the root mean square of the uncentered
coordinates and a learned scale.

The
\href{https://github.com/hermonai/inside-the-llm-engine/blob/astra-visual-rewrite/diagrams/transformer/layernorm-vs-rmsnorm.txt}{LayerNorm/RMSNorm
diagram} shows the difference in one glance. Removing the mean and
subtraction changes both mathematical invariances and execution. It is
not correct to explain RMSNorm as ``LayerNorm but faster'' without first
naming the missing operation and without measurement. The original
RMSNorm paper reports experimental efficiency and task results for its
studied systems; that evidence does not predict a universal
inference-kernel speed ratio on every engine.

\Appref{embeddings-norm} stops at this contrast. Batch normalization, group
normalization, partial RMSNorm, training gradients, and convergence
behavior are outside its inference boundary.

\section{Precision is part of the
operator}\label{precision-is-part-of-the-operator}

Real-number notation has unlimited range and precision. The Rust
operator does not. Transformer Primitives v1 makes one simple choice:

{\def\LTcaptype{none} % do not increment counter
\begin{longtable}[]{@{}ll@{}}
\toprule\noalign{}
Stage & Educational dtype \\
\midrule\noalign{}
\endhead
\bottomrule\noalign{}
\endlastfoot
input activation storage & \texttt{f32} \\
learned scale storage & \texttt{f32} \\
\texttt{x*x} product & \texttt{f32} \\
sum-of-squares accumulator & \texttt{f32} \\
mean, square root, reciprocal & \texttt{f32} \\
element-wise output product & \texttt{f32} \\
output storage & \texttt{f32} \\
\end{longtable}
}

The
\href{https://github.com/hermonai/inside-the-llm-engine/blob/astra-visual-rewrite/diagrams/transformer/normalization-precision-flow.txt}{precision-flow
diagram} places the finite-range gates along that path. There is no
implicit promotion to \texttt{f64} and no mixed-precision type
machinery.

Production execution can store activations or weights in FP16 or BF16
while using wider arithmetic for sensitive reductions. Meta's published
Llama 4 reference, for example, casts the normalization input to float
for square, mean, and reciprocal-square-root work, then casts the
normalized value back before applying the learned weight. This is one
source-verified model implementation, not a claim that every backend
follows the same sequence. Backend kernels can fuse stages while
preserving the model's accepted precision contract.

The reduction order matters too. In finite precision,

\begin{equation}
(a+b)+c
\end{equation}

can round differently from

\begin{equation}
a+(b+c).
\end{equation}

The scalar reference visits logical indices in increasing order. A
future SIMD tree, threaded partition, or GPU reduction can combine
partial sums in another order. Equivalent real-number equations
therefore need a numerical agreement policy: finite outputs are compared
with an absolute-plus-relative tolerance, while NaN or infinity cannot
pass merely because their neighboring values are close. Bit identity
remains appropriate for row selection and exact zero fixtures, not as a
universal cross-backend RMSNorm rule.

\section{Finite input can still
fail}\label{finite-input-can-still-fail}

\EngineFigure{ch07-range}{The magnitude sweep separates accepted finite reductions from square overflow and accumulated-sum overflow.}{fig:ch07-range}

Checking \texttt{x.is\_finite()} is necessary but insufficient. Squaring
expands the exponent. A finite \texttt{f32} around \texttt{1e20} has a
mathematical square around \texttt{1e40}, which exceeds the maximum
finite binary32 magnitude. \texttt{x*x} becomes infinity. Even when
every individual square is finite, their sum can overflow.

At the other end, very small squares can become subnormal or underflow
to zero. For tiny activations, epsilon can dominate the mean square. The
operator then behaves approximately like

\begin{equation}
y_i\approx\frac{x_iw_i}{\sqrt{\epsilon}},
\end{equation}

not like an exactly scale-invariant map.

The reference implementation remains the transparent naïve reduction. It
checks each square and the running sum and returns
\texttt{NonFiniteSquare} or \texttt{\seqsplit{NonFiniteReduction}}
instead of silently continuing. Detection does not expand the
algorithm's numerical range; it makes the limit observable.

Stable norm algorithms exist. Netlib's \texttt{SLASSQ}, for example,
represents a sum of squares using a scale and a scaled sum so
intermediate values can avoid many overflow and underflow cases. That is
valuable evidence and a possible future candidate. It does not belong in
the active \Appref{embeddings-norm} reference path: the more complex invariant would
obscure the direct equation, and there is no measured model workload
demanding substitution yet.

\begin{quote}
\textbf{PROVE IT}
\href{https://github.com/hermonai/inside-the-llm-engine/blob/astra-visual-rewrite/labs/lab-37-rmsnorm-magnitude-stress.md}{Lab
37} distinguishes finite input from finite intermediate arithmetic.
Remove the checks only in a disposable change and watch infinity turn a
normalization into a misleading zero or NaN.
\end{quote}

\section{Scale experiment: where epsilon becomes
visible}\label{scale-experiment-where-epsilon-becomes-visible}

Without epsilon and under exact real arithmetic, a positive scalar
\(\alpha\) cancels:

\begin{equation}
\frac{\alpha\mathbf{x}}
{\sqrt{\operatorname{mean}((\alpha\mathbf{x})^2)}}
=
\frac{\mathbf{x}}
{\sqrt{\operatorname{mean}(\mathbf{x}^2)}}.
\end{equation}

With fixed positive epsilon, the denominator becomes

\begin{equation}
\sqrt{\alpha^2m_2(\mathbf{x})+\epsilon},
\end{equation}

so exact cancellation no longer holds. It remains a good approximation
when \(\alpha^2m_2\gg\epsilon\).

The curve in the epsilon plate makes that condition visible. Keeping the
learned gain and epsilon fixed, exact real arithmetic gives the scalar
relation

\begin{equation}
\mathbf{y}(\alpha\mathbf{x})=c(\alpha)\mathbf{y}(\mathbf{x}),\qquad
c(\alpha)=\frac{\alpha\sqrt{m_2+\epsilon}}
{\sqrt{\alpha^2m_2+\epsilon}},\quad \alpha>0.
\end{equation}

Here \(c\) is the output scale relative to the unscaled input, not a
latency or accuracy score. For the hand vector, \(m_2=7.5\) and the two
denominator terms are equal at
\(\alpha=\sqrt{\epsilon/7.5}\approx0.00115\). Below that crossover,
epsilon dominates; above it, the curve approaches its nearly flat
regime. The horizontal axis is logarithmic so both behaviors remain
visible. This analytical curve complements the discrete executable
sweep, not a new timing benchmark or an assertion about trained
activation distributions.

The committed Rust example and independent Python oracle use the hand
vector, learned scale, and \texttt{epsilon=1e-5}. Relative to
\texttt{alpha=1}, the Python oracle records:

\Needspace{12\baselineskip}

{\def\LTcaptype{none} % do not increment counter
\begin{longtable}[]{@{}rr@{}}
\toprule\noalign{}
Positive scale \(\alpha\) & Maximum absolute output difference \\
\midrule\noalign{}
\endhead
\bottomrule\noalign{}
\endlastfoot
\texttt{1e-8} & \texttt{2.1908698} \\
\texttt{0.1} & \texttt{0.000144584} \\
\texttt{1} & \texttt{0} \\
\texttt{10} & \texttt{0.000001446} \\
\texttt{100} & \texttt{0.000001460} \\
\end{longtable}
}

The tiny-scale counterexample matters more than the close large-scale
rows. A blanket sentence ``RMSNorm is scale invariant'' would be false
for the operator we actually execute.
\href{https://github.com/hermonai/inside-the-llm-engine/blob/astra-visual-rewrite/labs/lab-36-rmsnorm-scale-experiment.md}{Lab
36} reproduces the sweep and asks where the epsilon crossover lies.

The magnitude sweep adds another view:

\Needspace{12\baselineskip}

{\def\LTcaptype{none} % do not increment counter
\begin{longtable}[]{@{}rl@{}}
\toprule\noalign{}
Input magnitude & Teaching \texttt{f32} observation with alternating
signs \\
\midrule\noalign{}
\endhead
\bottomrule\noalign{}
\endlastfoot
\texttt{1e-20} & finite subnormal square; output scale dominated by
epsilon \\
\texttt{1e-10} & finite square; epsilon still dominates \\
\texttt{1} & output approximately \texttt{±0.999995} \\
\texttt{1e10} & finite square and output \texttt{±1} at printed
precision \\
\texttt{1e20} & typed square-overflow error \\
\end{longtable}
}

Four values of \texttt{1e19} separately square to finite values but
overflow their running \texttt{f32} sum, exercising a different error.
These are bounded numerical experiments, not a claim that real trained
activations span this entire range.

\section{Work and memory behavior}\label{work-and-memory-behavior}

\EngineFigure{ch07-cost}{Four F32 payload streams explain an analytical 16D-byte RMSNorm cost without asserting measured memory traffic.}{fig:ch07-cost}

Embedding lookup and RMSNorm differ from \Appref{matmul}'s matrix products.

A single embedding lookup selects one row, checks addresses, reads \(D\)
values, and under our policy writes \(D\) values. It performs almost no
floating-point arithmetic. Vocabulary size controls parameter-table
capacity and ID bounds, but a single valid lookup does not scan all
\(V_{\mathrm{vocab}}\) rows.

The RMSNorm reference performs a reduction and an element-wise pass.
Under a simple cold-payload \texttt{f32} model it:

\begin{enumerate}
\def\labelenumi{\arabic{enumi}.}
\tightlist
\item
  reads input once for the reduction: \(4D\) bytes;
\item
  reads input again for scaling: \(4D\) bytes;
\item
  reads learned weight: \(4D\) bytes;
\item
  writes output: \(4D\) bytes.
\end{enumerate}

Thus

\begin{equation}
Q_{\mathrm{RMSNorm}}\approx16D\quad[\mathrm{bytes}].
\end{equation}

This estimate excludes tensor metadata, output allocation overhead,
cache-line effects, write allocation, and reuse from nearby operations.
The computation does roughly linear work in \(D\), including a square
and sum per input and two multiplications per output, plus division,
square root, and reciprocal work. Assigning one universal FLOP count to
square root or reciprocal is not useful here.

The structural conclusion is narrower: this operation exposes much less
data reuse than a large GEMM, in which loaded matrix blocks can feed
many output cells. Its arithmetic intensity is low under the stated
payload boundary, so optimization questions will often concern data
movement, reduction strategy, and fusion. ``Often'' is architectural
reasoning, not a measurement of the current example.

No timing benchmark is published for \Appref{embeddings-norm}. A nanosecond result for
one tiny vector would mostly expose call, allocation, compiler, and
timer details. The scale and magnitude experiments answer actual
correctness questions. A later optimization should first define the
candidate, workload, cache state, repetitions, hardware, compiler, and
equivalence gate before measuring speed.

\section{Embedding is not output
projection}\label{embedding-is-not-output-projection}

ENGINE-1 contains both an embedding table and an output weight matrix
with shape \texttt{{[}V,D{]}}. Equal shape does not make the operations
equal.

Embedding maps a token identity to a model-space row:

\begin{equation}
t\longmapsto\mathbf{E}_{t,:}\in\mathbb{R}^{D}.
\end{equation}

Output projection maps a hidden vector to one score per vocabulary item:

\begin{equation}
\mathbf{z}=\mathbf{W}_{\mathrm{out}}\mathbf{x}+\mathbf{b},
\qquad
\mathbf{W}_{\mathrm{out}}\in
\mathbb{R}^{V_{\mathrm{vocab}}\times D},
\quad
\mathbf{z}\in\mathbb{R}^{V_{\mathrm{vocab}}}.
\end{equation}

Lookup reads one selected row. Projection uses every output row in a
GEMV. The
\href{https://github.com/hermonai/inside-the-llm-engine/blob/astra-visual-rewrite/diagrams/transformer/embedding-vs-output-projection.txt}{comparison
diagram} shows where the two \texttt{{[}V,D{]}} tensors act.

Some architectures tie embedding and output-projection weights, sharing
parameter storage. That is a model architecture decision, not a
consequence of shape. ENGINE-1 keeps distinct owners and \Appref{embeddings-norm} does
not redesign it.

\section{Integration without rewriting
history}\label{integration-without-rewriting-history}

Before this chapter,
\texttt{\seqsplit{TinyLanguageModel::embedding\_row}} manually checked
the token and copied each tensor element. It now calls
\texttt{\seqsplit{embedding\_lookup\_reference}}, converts the owned
tensor into its existing hidden
\texttt{Vec\textless{}f32\textgreater{}}, and maps an invalid token back
to the established \texttt{\seqsplit{ModelError::TokenOutOfRange}}
variant. The output GEMV and bias addition are unchanged.

RMSNorm is \emph{not} inserted into ENGINE-1. That tiny model was
defined without a normalization weight or normalization step. Adding one
merely to claim integration would change its logits and invalidate the
historical hand fixture. The new operator is a tested primitive for the
Transformer path that later chapters will compose honestly.

The known regression remains:

\begin{verbatim}
input token:       like
embedding/hidden: [1.0,-0.5,2.0]
logits:            [-0.7,0.1,0.4,2.2]
greedy output:     Rust, then EOS
\end{verbatim}

This is a useful architectural lesson: extracting an operation should
preserve the model's semantics. Adding a new primitive to a library does
not authorize changing old model graphs.

\section{Independent correctness}\label{independent-correctness}

The visual edition adds a second, explicitly connected evidence path:

\begin{verbatim}
python3 scripts/check-normalization-visual-parity.py
\end{verbatim}

An input-only JSON fixture supplies the table, token sequence, gain and
epsilon. The Rust example calls the real operators; it does not print a
stored expected answer. The independent Python computation supplies both
a wider mathematical answer and an explicitly rounded F32 trace. The
gate compares the entire output, selected addresses, squares and partial
sums, then checks that all ten canonical scenes are embedded once. Five
additional Rust tests cover the connecting fixture and its ownership and
range claims. Existing oracle and stress tests remain independent
regressions, not replaced by screenshots.

The original operator milestone added 30 deterministic tests. Embedding
coverage includes first, middle, and last rows; dimensions one and
greater; wrong rank; empty axes; out-of-range IDs; strided and
zero-stride views; sequence order and repetition; empty token sequence;
and ownership isolation.

RMSNorm coverage includes dimension one; the hand vector; mixed signs;
zero and uniform vectors; non-unit weights; strided and zero-stride
views; wrong ranks; length mismatch; empty dimension; invalid epsilon;
non-finite operands; large and small magnitudes; square, reduction, and
output overflow; and approximate positive-scale invariance where epsilon
is negligible.

The Python oracle is independently expressed. It uses ordinary Python
lists, \texttt{math.fsum}, and Python's wider floating-point evaluation
for the primary mathematical result. A separate \texttt{struct}-based
helper rounds every stress operation to IEEE binary32 so it can classify
the Rust reference's finite range. It does not import Rust, NumPy, or
another tensor framework.

Independent expression reduces correlated mistakes: the Rust operator
follows checked tensor strides, while Python works from mathematical
lists. It does not prove every possible vector. Boundary tests,
deterministic fixtures, and the unchanged previous oracles remain part
of the gate.

\begin{quote}
\textbf{PROVE IT}
\href{https://github.com/hermonai/inside-the-llm-engine/blob/astra-visual-rewrite/labs/lab-38-rust-python-rmsnorm.md}{Lab
38} compares exact lookup results and tolerant RMSNorm results.
Deliberately permute learned weights; an RMS-only scalar check will miss
the bug, while full-vector comparison will catch it.
\end{quote}

\section{Inside Hermon}\label{inside-hermon-3}

\EngineFigure{ch07-source}{Source-verified default and preview routes separate graph execution from the visible Rust normalization loop.}{fig:ch07-source}

The following paths were rechecked at Hermon commit \texttt{2a3fd521} on
2026-09-10; the full pin and original historical observations are
retained in the
\href{https://github.com/hermonai/inside-the-llm-engine/blob/astra-visual-rewrite/research/astra/chapter07-regeneration.md}{visual
review record}. Status labels matter because Hermon has a release path
and a separately gated engine preview.

\begin{quote}
\textbf{INSIDE HERMON --- CURRENT} With
\texttt{\seqsplit{HERMON\_RUNTIME\_MODE}} unset, \texttt{dispatch.rs}
selects \texttt{Batched}. That path constructs the llama.cpp-backed
batched runtime. Embedding lookup, normalization, and the rest of the
default model graph therefore execute in the pinned llama.cpp/GGML
machinery rather than the Rust reference loop in \texttt{paged.rs}.
\end{quote}

\begin{quote}
\textbf{INSIDE HERMON --- PREVIEW} Explicit
\texttt{\seqsplit{HERMON\_RUNTIME\_MODE=paged}} selection is accompanied
by source-level preview warnings and additional gates. Its
\texttt{\seqsplit{GgufLlamaForward}} reads token embedding rows through
a typed llama.cpp bridge, loads normalization vectors, validates model
epsilon, and runs a visible Rust \texttt{rms\_norm} function for its
Hermon-owned forward path.
\end{quote}

\begin{quote}
\textbf{INSIDE HERMON --- LIBRARY}
\texttt{\seqsplit{hermon-gguf::model\_shape}} verifies that
\texttt{\seqsplit{token\_embd.weight}} begins with GGML dimension
\texttt{\seqsplit{embedding\_length}} and derives vocabulary size from
the next dimension.
\texttt{\seqsplit{hermon-llamacpp::tensor\_row\_f32}} validates a
requested row and materializes it as a Rust
\texttt{Vec\textless{}f32\textgreater{}}, converting supported
host-resident packed storage through the C++ bridge. These library
abilities are consumed by the paged preview; their existence does not
make that path the default.
\end{quote}

The preview Rust function computes \texttt{sum(x*x)} in \texttt{f32},
divides by length, adds epsilon inside \texttt{sqrt}, takes the
reciprocal, and multiplies input and gain. Its debug assertions rely on
model construction to establish equal slice lengths. The teaching API is
more general and defensive because it is a public checked operator over
arbitrary valid views.

The
\href{https://github.com/hermonai/inside-the-llm-engine/blob/astra-visual-rewrite/diagrams/transformer/hermon-llamacpp-normalization-path.txt}{Hermon
mapping diagram} keeps CURRENT, PREVIEW, and LIBRARY boundaries on the
same page. It also prevents a common source-reading error: finding
native code does not prove the default API request reaches it.

\section{Inside llama.cpp and
GGML}\label{inside-llama.cpp-and-ggml-1}

Hermon's submodule pins llama.cpp/GGML commit
\texttt{\seqsplit{389ff61d77b5c71cec0cf92fe4e5d01ace80b797}}. At that
revision the architecture has several layers.

The llama graph builder creates token embeddings with
\texttt{ggml\_get\_rows} over the token-embedding tensor and input-token
tensor. Its \texttt{build\_norm} function maps the model's RMS
normalization type to \texttt{ggml\_rms\_norm}, passing
\texttt{f\_norm\_rms\_eps}. Learned multiplication is represented by
subsequent graph work and can participate in fused execution.

\texttt{ggml\_get\_rows} first creates an operator node and result
tensor. The result's fastest physical dimension is the source row width;
remaining dimensions track the index tensor. At CPU execution, backend
dispatch selects by source storage type:

\begin{itemize}
\tightlist
\item
  F32 rows are copied;
\item
  F16 and BF16 rows are converted to F32;
\item
  supported quantized rows are dequantized to F32.
\end{itemize}

Thus a ``row lookup'' can include conversion and materialization in a
production engine. The graph knows the semantic operation; the backend
kernel knows the physical type.

The inspected CPU F32 RMSNorm kernel requires the first tensor dimension
to be contiguous. It distributes outer rows across threads and visits
each contiguous inner row. At this commit \texttt{ggml-cpu/vec.h}
defines the \texttt{ggml\_float} accumulator as \texttt{double}, but
\texttt{x{[}i{]}*x{[}i{]}} is formed from two F32 operands before that
product is cast into the wider sum. It computes
\texttt{mean\ =\ sum\ /\ row\_width}, then
\texttt{scale\ =\ 1/sqrtf(mean\ +\ eps)}. The plain path copies and
scales the row. A fused RMSNorm-plus-multiply path applies the learned
weight during the output loop.

GGML stores epsilon in operator parameters and asserts a nonnegative
value in that kernel. It asserts row indices and layout invariants after
graph and model construction. Those assertions are appropriate inside
its validated execution graph; the public educational functions instead
return typed errors for ordinary malformed calls.

The source flow is:

The model graph declares RMS normalization and learned scaling. The CPU execution path reads epsilon, accumulates squares, computes a scale, then writes the normalized values; fusion can combine the learned multiply with that write.

This chapter describes the pinned CPU path, not every backend.
Accelerator kernels can partition reductions, fuse operations, and
impose other layout or dtype contracts. Source at one commit also cannot
justify ``llama.cpp always'' claims across revisions.

\section{Educational and production
boundaries}\label{educational-and-production-boundaries}

The teaching path is intentionally narrow:

Immutable strided F32 inputs enter a checked two-pass reference and produce a fresh canonical F32 owner. No optimized normalization candidate is claimed.

A production engine may need lower-precision activation storage, packed
or quantized parameter rows, wider accumulation, vector or thread
reductions, device-local addresses, graph scheduling, batch dimensions,
and fused normalization with neighboring work. Each complication changes
a contract: conversion precision, work partition, ownership, layout, or
failure placement.

The reference remains valuable precisely because it does not absorb all
those concerns. A future optimized candidate can be compared against one
named semantic path. The permanent substitution order remains:

First establish scalar semantics and an independent oracle. Next check any proposed replacement against them. Measure it only after correctness passes.

\Appref{embeddings-norm} has no candidate because correctness infrastructure should
precede optimization pressure. That is not unfinished work. It is a
deliberate engine state with an inspectable active implementation.

\section{Common mistakes}\label{common-mistakes-5}

Several short statements conceal major errors:

\begin{itemize}
\tightlist
\item
  \textbf{``The token ID is fed into the neural network as a number.''}
  The ID selects a learned row; its scalar magnitude is not an
  activation feature.
\item
  \textbf{``Embedding is a matrix multiply.''} One-hot multiplication is
  an equivalence, while the implemented operation is indexed row
  materialization.
\item
  \textbf{``Shape \texttt{{[}V,D{]}} guarantees contiguous rows.''} Only
  shape plus strides and storage metadata establish layout.
\item
  \textbf{``Returning a view is always faster.''} It avoids one copy but
  constrains lifetime and preserves aliasing; safety and later mutation
  requirements decide the interface.
\item
  \textbf{``RMSNorm makes every component unit length.''} It uses one
  vector-wide RMS factor and then learned per-coordinate scale.
  Individual coordinates need not have magnitude one.
\item
  \textbf{``RMSNorm subtracts the mean.''} That describes the centering
  found in LayerNorm, not this operator.
\item
  \textbf{``Epsilon can go anywhere near the denominator.''} Placement
  defines the function.
\item
  \textbf{``Finite input means finite squares.''} A finite \texttt{f32}
  can overflow when squared, and finite squares can overflow their sum.
\item
  \textbf{``Equivalent reductions must match bit-for-bit.''} Different
  valid reduction trees can round differently; equivalence needs a
  justified tolerance and finite-value gate.
\item
  \textbf{``A Rust RMSNorm function proves Hermon uses it in
  production.''} Runtime selection proves the current default path; at
  the inspected commit it is llama.cpp-backed batched execution.
\end{itemize}

\section{Transformer Primitives v1}\label{transformer-primitives-v1}

\Appref{embeddings-norm} completes a named extension of ENGINE-2 rather than inventing
ENGINE-3. \textbf{Transformer Primitives v1} means the repository can
now execute:

\begin{itemize}
\tightlist
\item
  checked single-token embedding
  \texttt{{[}V,D{]}\ +\ TokenId\ -\textgreater{}\ {[}D{]}};
\item
  checked sequence embedding
  \texttt{{[}V,D{]}\ +\ {[}T{]}\ -\textgreater{}\ {[}T,D{]}};
\item
  explicit parameter-to-activation copy ownership;
\item
  checked scalar RMSNorm
  \texttt{{[}D{]}\ +\ {[}D{]}\ +\ epsilon\ -\textgreater{}\ {[}D{]}};
\item
  valid strided and zero-stride reference layouts;
\item
  typed shape, epsilon, ID, and finite-range failures;
\item
  independent Python and Rust numerical verification.
\end{itemize}

ENGINE-2's dot, GEMV, reference GEMM, and blocked GEMM remain intact.
The
\href{https://github.com/hermonai/inside-the-llm-engine/blob/astra-visual-rewrite/diagrams/transformer/chapter07-engine-architecture.txt}{milestone
architecture} places all three ingredients---embedding, normalization,
and linear algebra---on the checked tensor substrate without composing a
fake Transformer layer.

The labs form a complete evidence path:

\begin{itemize}
\tightlist
\item
\href{https://github.com/hermonai/inside-the-llm-engine/blob/astra-visual-rewrite/labs/README.md}{Labs
  30--32} follow embedding offsets, checked lookup, and ownership;
\item
\href{https://github.com/hermonai/inside-the-llm-engine/blob/astra-visual-rewrite/labs/README.md}{Labs
  33--35} derive RMSNorm and break its epsilon contract;
\item
\href{https://github.com/hermonai/inside-the-llm-engine/blob/astra-visual-rewrite/labs/README.md}{Labs
  36--37} expose scale and finite-range behavior;
\item
  \href{https://github.com/hermonai/inside-the-llm-engine/blob/astra-visual-rewrite/labs/lab-38-rust-python-rmsnorm.md}{Lab
  38} closes the independent oracle gate.
\end{itemize}

\section{What we can now account
for}\label{what-we-can-now-account-for}

For embedding lookup, we can name the long-lived parameter owner, both
tensor axes, the selected logical coordinates, canonical and strided
offsets, copied bytes, activation owner, valid ID interval, and every
empty-axis decision.

For RMSNorm, we can name input and learned-scale shapes, the reduced
dimension, epsilon convention, storage and reduction dtype, loop order,
two-pass traffic, output owner, zero-vector result, invalid
configuration, overflow boundary, oracle, and tolerance policy.

That is the level at which later kernels must be understood. The
equation, memory, ownership, and production mapping agree. Short code is
no excuse for a vague system contract.

\section{The boundary to Chapter~\ref{app:qkv}}\label{the-boundary-to-chapter-8}

We now have model-space vectors, normalization, and matrix
multiplication. How does a Transformer turn one residual vector into
three different representations that determine what information it asks
for, what information it offers, and what information is ultimately
transported?

That is the bounded question for \textbf{Chapter~\ref{app:qkv} --- Queries, Keys,
and Values}. \Appref{qkv} may build the three linear projections and their
shape/ownership contracts. It must not smuggle in attention, position
rotation, or KV caching: those remain separate later chapters.

\section{Primary references}\label{primary-references-2}

\begin{itemize}
\tightlist
\item
  Biao Zhang and Rico Sennrich,
  \href{https://arxiv.org/abs/1910.07467}{``Root Mean Square Layer
  Normalization''}, 2019;
  \href{https://papers.neurips.cc/paper_files/paper/2019/file/1e8a19426224ca89e83cef47f1e7f53b-Paper.pdf}{NeurIPS
  proceedings version}.
\item
  Jimmy Lei Ba, Jamie Ryan Kiros, and Geoffrey E. Hinton,
  \href{https://arxiv.org/abs/1607.06450}{``Layer Normalization''},
  2016.
\item
  PyTorch,
  \href{https://docs.pytorch.org/docs/2.14/generated/torch.nn.RMSNorm.html}{\texttt{\seqsplit{torch.nn.RMSNorm}}},
  rechecked 2026-09-10, and
  \href{https://docs.pytorch.org/docs/stable/generated/torch.nn.Embedding.html}{\texttt{\seqsplit{torch.nn.Embedding}}},
  accessed 2026-09-03.
\item
  Meta,
  \href{https://github.com/meta-llama/llama-models/blob/main/models/llama4/model.py}{Llama
  4 reference \texttt{model.py}}, accessed 2026-09-03.
\item
  Netlib LAPACK,
  \href{https://www.netlib.org/lapack/explore-html/d8/d76/group__lassq_ga0596b4bfa745d0d1c5817d4790921cda.html}{\texttt{SLASSQ}:
  scaled sum of squares}, accessed 2026-09-03.
\item
  Rust standard library,
  \href{https://doc.rust-lang.org/std/primitive.f32.html}{\texttt{f32}},
  accessed 2026-09-03.
\item
  Hermon source at \texttt{2a3fd521} and its pinned llama.cpp/GGML
  source at \texttt{389ff61d}, rechecked 2026-09-10; full hashes and
  classifications are in the
  \href{https://github.com/hermonai/inside-the-llm-engine/blob/astra-visual-rewrite/research/astra/chapter07-regeneration.md}{visual
  review}. The
  \href{https://github.com/hermonai/inside-the-llm-engine/blob/astra-visual-rewrite/research/part-02/chapter-07-embeddings-and-normalization.md}{original
  research note} preserves the 2026-09-03 historical inspection.
\end{itemize}


\section{Worked synthesis: normalization does not erase everything}
\subsection{Compute every term}
\textbf{Problem.} Let $x=(3,4)$, $\gamma=(1,2)$ and $\epsilon=.5$.
Compute RMSNorm in real arithmetic.

\textbf{Solution.} The mean square is $(9+16)/2=12.5$; the denominator is
$\sqrt{13}$. Therefore
\begin{equation}
y=\left(\frac3{\sqrt{13}},\frac8{\sqrt{13}}\right).
\end{equation}
The second learned scale is applied after the common normalization. The
mean of $y$ is not zero, and its RMS need not be one because $\gamma$ and
$\epsilon$ change the result.

\subsection{Derive the limit of scale invariance}
\textbf{Problem.} Compare the denominator for $ax$ with that for $x$.

\textbf{Solution.} If $s=\frac1D\sum_i x_i^2$, the denominator is
$\sqrt{a^2s+\epsilon}$. For positive $a$,
$ax/\sqrt{a^2s+\epsilon}=x/\sqrt{s+\epsilon/a^2}$.
It equals the unscaled expression only under special conditions, including
$\epsilon=0$. For negative $a$, the sign also reverses. Numerical overflow in
the square is a further implementation concern, not addressed by this
real-arithmetic manipulation.

\subsection{Preserve the earlier fixture}
\textbf{Problem.} Why not silently insert RMSNorm into the \Appref{smallest-model} model
as soon as the operator exists?

\textbf{Solution.} Doing so changes its logits and the meaning of earlier
regression examples. A new composed model must explicitly declare the new
operation and receive its own expected outputs. Component availability is
not permission to change the mathematical model behind an established test.

